% !TEX root = ../../main.tex
\chapter*{抽象代数名词索引}
\phantomsection
\addcontentsline{toc}{chapter}{抽象代数名词索引}
\markboth{抽象代数名词索引}{抽象代数名词索引}

\begingroup
\setlength{\tabcolsep}{8pt}
\renewcommand{\arraystretch}{1.25}
\begin{longtable}{@{}>{\raggedright\arraybackslash}p{0.47\linewidth}>{\raggedright\arraybackslash}p{0.47\linewidth}@{}}
\hline
\endfirsthead
\hline
\endhead
\hline
\endfoot
\hline
\endlastfoot
\hyperref[p2:concept:c3-divisibility]{倍式} & \hyperref[p2:concept:c2-l405]{被移动的点}\\
\hyperref[p2:concept:c1-l248]{变换群} & \hyperref[p2:concept:c2-l1600]{不可分解群}\\
\hyperref[p2:concept:c3-irreducible]{不可约多项式} & \hyperref[p2:concept:c1-l210]{不相交轮换}\\
\hyperref[p2:concept:c3-functions]{常函数} & \hyperref[p2:concept:c3-constant-subfield]{常函数子域}\\
\hyperref[p2:concept:c3-evaluation]{超越元} & \hyperref[p2:concept:c1-l954]{乘法单位元}\\
\hyperref[p2:concept:c3-multiplicative-set]{乘性子集（乘法子集）} & \hyperref[p2:concept:c2-l1035]{传递作用}\\
\hyperref[p2:concept:c3-evaluation]{代入同态} & \hyperref[p2:concept:c3-algebraic-relation]{代数关系}\\
\hyperref[p2:concept:c3-finitely-generated-ring]{代数无关} & \hyperref[p2:concept:c3-finitely-generated-ring]{代数相关}\\
\hyperref[p2:concept:c3-evaluation]{代数元} & \hyperref[p2:concept:c1-l19]{代数运算（二元运算）}\\
\hyperref[p2:concept:c2-l332]{单群} & \hyperref[p2:concept:c1-l956]{单位（可逆元）}\\
\hyperref[p2:concept:c1-l963]{单位群} & \hyperref[p2:concept:c1-l40]{单位元}\\
\hyperref[p2:concept:c3-multivariate-ring]{单项式} & \hyperref[p2:concept:c2-l504]{导出列}\\
\hyperref[p2:concept:c2-l504]{导出子群} & \hyperref[p2:concept:c2-l1004]{等价的作用}\\
\hyperref[p2:concept:c1-l193]{$n$元对称群} & \hyperref[p2:concept:c1-l210]{对换}\\
\hyperref[p2:concept:c3-degree]{多项式的次数} & \hyperref[p2:concept:c3-evaluation]{多项式的根}\\
\hyperref[p2:concept:c3-divisibility]{多项式的整除} & \hyperref[p2:concept:c3-polynomial-functions]{多项式函数}\\
\hyperref[p2:concept:c3-polynomial-functions]{多项式函数环} & \hyperref[p2:concept:c3-multivariate-functions]{多元多项式函数环}\\
\hyperref[p2:concept:c3-multivariate-ring]{多元多项式环} & \hyperref[p2:concept:c1-l164]{二面体群}\\
\hyperref[p2:concept:c3-nilpotent]{非零幂元} & \hyperref[p2:concept:c3-finitely-generated-ring]{非平凡代数关系}\\
\hyperref[p2:concept:c1-l253]{非平凡子群} & \hyperref[p2:concept:c3-fraction-ring]{分式环}\\
\hyperref[p2:concept:c2-l1151]{共轭类} & \hyperref[p2:concept:c2-l978]{共轭元素}\\
\hyperref[p2:concept:c2-l978]{共轭子群} & \hyperref[p2:concept:c2-l971]{共轭作用}\\
\hyperref[p2:concept:c2-l1031]{轨道} & \hyperref[p2:concept:c3-multivariate-functions]{函数环}\\
\hyperref[p2:concept:c3-ideal-product]{互素理想} & \hyperref[p2:concept:c1-l781]{环}\\
\hyperref[p2:concept:c3-antiisomorphism]{环的反同构} & \hyperref[p2:concept:c3-antiisomorphism]{环的反同态}\\
\hyperref[p2:concept:c3-antiisomorphism]{环的反自同构} & \hyperref[p2:concept:c1-l785]{环的零元}\\
\hyperref[p2:concept:c1-l1157]{环的满同态} & \hyperref[p2:concept:c3-ring-internal-sum]{环的内直和}\\
\hyperref[p2:concept:c3-ring-sum]{环的直和} & \hyperref[p2:concept:c1-l1328]{环的自然同态}\\
\hyperref[p2:concept:c1-l923]{环同构} & \hyperref[p2:concept:c1-l1152]{环同态}\\
\hyperref[p2:concept:c1-l1156]{环同态的核} & \hyperref[p2:concept:c1-l1156]{环同态的像}\\
\hyperref[p2:concept:c1-l830]{环中的整数倍} & \hyperref[p2:concept:c2-l173]{换位子}\\
\hyperref[p2:concept:c2-l173]{换位子群} & \hyperref[p2:concept:c3-prime-maximal]{极大理想}\\
\hyperref[p2:concept:c2-l617]{极大正规子群} & \hyperref[p2:concept:c3-minimal-polynomial]{极小多项式}\\
\hyperref[p2:concept:c1-l106]{加法记号下的整数倍} & \hyperref[p2:concept:c1-l785]{加法逆元}\\
\hyperref[p2:concept:c1-l290]{交错群} & \hyperref[p2:concept:c1-l958]{交换环}\\
\hyperref[p2:concept:c1-l37]{交换律} & \hyperref[p2:concept:c1-l60]{交换群（阿贝尔群）}\\
\hyperref[p2:concept:c1-l762]{交换图} & \hyperref[p2:concept:c1-l34]{结合律}\\
\hyperref[p2:concept:c2-l522]{可解群} & \hyperref[p2:concept:c3-irreducible]{可约多项式}\\
\hyperref[p2:concept:c3-extension-ring]{扩环} & \hyperref[p2:concept:c3-subring-sum]{理想的和}\\
\hyperref[p2:concept:c3-ideal-product]{理想的积} & \hyperref[p2:concept:c3-zero-function]{零函数}\\
\hyperref[p2:concept:c1-l1261]{零环} & \hyperref[p2:concept:c3-nilpotent-ideal]{零幂理想}\\
\hyperref[p2:concept:c3-nilpotent]{零幂元} & \hyperref[p2:concept:c1-l1157]{零同态}\\
\hyperref[p2:concept:c1-l989]{零因子} & \hyperref[p2:concept:c1-l209]{轮换}\\
\hyperref[p2:concept:c2-l372]{轮换类型} & \hyperref[p2:concept:c3-congruence]{模理想同余}\\
\hyperref[p2:concept:c3-primitive-root]{模$p$的原根} & \hyperref[p2:concept:c2-l1542]{内直和}\\
\hyperref[p2:concept:c2-l790]{内自同构} & \hyperref[p2:concept:c2-l799]{内自同构群}\\
\hyperref[p2:concept:c1-l42]{逆元} & \hyperref[p2:concept:c1-l1356]{欧拉函数}\\
\hyperref[p2:concept:c1-l289]{偶置换} & \hyperref[p2:concept:c1-l312]{陪集代表}\\
\hyperref[p2:concept:c1-l1203]{平凡理想} & \hyperref[p2:concept:c1-l252]{平凡子群}\\
\hyperref[p2:concept:c2-l986]{齐性空间} & \hyperref[p2:concept:c1-l289]{奇置换}\\
\hyperref[p2:concept:c1-l192]{全变换群} & \hyperref[p2:concept:c1-l52]{群}\\
\hyperref[p2:concept:c1-l149]{群的阶} & \hyperref[p2:concept:c1-l487]{群的满同态}\\
\hyperref[p2:concept:c1-l487]{群的嵌入} & \hyperref[p2:concept:c2-l808]{群的中心}\\
\hyperref[p2:concept:c1-l727]{群的自然同态} & \hyperref[p2:concept:c2-l924]{群的左作用}\\
\hyperref[p2:concept:c1-l388]{群同构} & \hyperref[p2:concept:c1-l482]{群同态}\\
\hyperref[p2:concept:c1-l483]{群同态的核} & \hyperref[p2:concept:c1-l483]{群同态的像}\\
\hyperref[p2:concept:c1-l1252]{商环} & \hyperref[p2:concept:c1-l681]{商群}\\
\hyperref[p2:concept:c3-fraction-field]{商域} & \hyperref[p2:concept:c2-l146]{生成元（生成元组）}\\
\hyperref[p2:concept:c3-degree]{首项系数} & \hyperref[p2:concept:c3-degree]{首一多项式}\\
\hyperref[p2:concept:c1-l1201]{双边理想（理想）} & \hyperref[p2:concept:c1-l1092]{四元数体}\\
\hyperref[p2:concept:c3-prime-maximal]{素理想} & \hyperref[p2:concept:c1-l1517]{素域}\\
\hyperref[p2:concept:c2-l891]{特征子群} & \hyperref[p2:concept:c1-l1037]{体（除环）}\\
\hyperref[p2:concept:c1-l388]{同构的群} & \hyperref[p2:concept:c1-l156]{图形的对称群}\\
\hyperref[p2:concept:c2-l1501]{外直和} & \hyperref[p2:concept:c2-l845]{外自同构群}\\
\hyperref[p2:concept:c2-l885]{完全群} & \hyperref[p2:concept:c3-indeterminate]{未定元（不定元）}\\
\hyperref[p2:concept:c2-l1046]{稳定子群} & \hyperref[p2:concept:c1-l149]{无限群}\\
\hyperref[p2:concept:c2-l1287]{西罗$p$子群} & \hyperref[p2:concept:c3-formal-series]{形式幂级数}\\
\hyperref[p2:concept:c2-l204]{循环群} & \hyperref[p2:concept:c1-l954]{幺环}\\
\hyperref[p2:concept:c3-polynomial-ring]{一元多项式环} & \hyperref[p2:concept:c3-formal-series]{一元形式幂级数环}\\
\hyperref[p2:concept:c3-nilradical]{诣零根} & \hyperref[p2:concept:c3-divisibility]{因式}\\
\hyperref[p2:concept:c1-l353]{由单个元素生成的子群} & \hyperref[p2:concept:c3-generated-ideal]{由集合生成的理想}\\
\hyperref[p2:concept:c2-l142]{由集合生成的子群} & \hyperref[p2:concept:c3-generated-subring]{由元素生成的子环}\\
\hyperref[p2:concept:c2-l1505]{有限个群的直积} & \hyperref[p2:concept:c1-l149]{有限群}\\
\hyperref[p2:concept:c3-finitely-generated-ring]{有限生成环} & \hyperref[p2:concept:c3-generated-ideal]{有限生成理想}\\
\hyperref[p2:concept:c2-l146]{有限生成群} & \hyperref[p2:concept:c2-l1111]{有限$p$群}\\
\hyperref[p2:concept:c1-l1202]{右理想} & \hyperref[p2:concept:c1-l989]{右零因子}\\
\hyperref[p2:concept:c1-l311]{右陪集} & \hyperref[p2:concept:c1-l468]{右正则表示}\\
\hyperref[p2:concept:c1-l1035]{域} & \hyperref[p2:concept:c1-l1442]{域的特征}\\
\hyperref[p2:concept:c3-functions]{域上的函数} & \hyperref[p2:concept:c1-l354]{元素的阶}\\
\hyperref[p2:concept:c1-l488]{元素的完全反像} & \hyperref[p2:concept:c1-l105]{元素的整数次幂}\\
\hyperref[p2:concept:c1-l990]{整环} & \hyperref[p2:concept:c2-l1208]{正规化子}\\
\hyperref[p2:concept:c1-l574]{正规子群} & \hyperref[p2:concept:c1-l830]{正整数幂}\\
\hyperref[p2:concept:c1-l194]{置换} & \hyperref[p2:concept:c2-l405]{置换的不动点}\\
\hyperref[p2:concept:c1-l284]{置换的符号} & \hyperref[p2:concept:c2-l1162]{中心化子}\\
\hyperref[p2:concept:c2-l1001]{忠实（如实）作用} & \hyperref[p2:concept:c3-generated-ideal]{主理想}\\
\hyperref[p2:concept:c3-pid]{主理想整环} & \hyperref[p2:concept:c3-functions]{自变量函数}\\
\hyperref[p2:concept:c2-l718]{自同构} & \hyperref[p2:concept:c2-l722]{自同构群}\\
\hyperref[p2:concept:c1-l887]{子环} & \hyperref[p2:concept:c3-subring-sum]{子环的和}\\
\hyperref[p2:concept:c1-l613]{子集的乘积} & \hyperref[p2:concept:c1-l613]{子集的逆集}\\
\hyperref[p2:concept:c1-l247]{子群} & \hyperref[p2:concept:c2-l12]{子群的完全反像}\\
\hyperref[p2:concept:c1-l312]{子群的指数} & \hyperref[p2:concept:c1-l1039]{子域}\\
\hyperref[p2:concept:c1-l55]{左单位元} & \hyperref[p2:concept:c1-l1202]{左理想}\\
\hyperref[p2:concept:c1-l988]{左零因子} & \hyperref[p2:concept:c1-l56]{左逆元}\\
\hyperref[p2:concept:c1-l311]{左陪集} & \hyperref[p2:concept:c2-l963]{左平移作用}\\
\hyperref[p2:concept:c1-l468]{左正则表示} & \hyperref[p2:concept:c2-l1035]{作用的不动点及全体不动点}\\
\end{longtable}
\endgroup
